% Vietnamese translation for OI Public Library
% Source-File: IOI中国国家候选队论文/2003/张宁/张宁.ppt
\documentclass[11pt]{article}
\usepackage{oipl-vn}

\title{Bản trình chiếu: Nghiên cứu bài toán đoán số}
\author{Zhang Ning}
\date{2003}

\begin{document}
\maketitle

\origtitle{Nghiên cứu bài toán đoán số}
\sourcepath{IOI China candidate team papers / 2003 / Zhang Ning / Zhang Ning.ppt}

\section{Mục tiêu của slide}

Bộ slide trình bày lại bài toán đoán số, một ví dụ điển hình của suy luận
``lồng tư duy''. Nhiều công thức trong PowerPoint là đối tượng nhúng nên phần
chữ trích được khá ít; bản ghi chú này dựa trên bài viết cùng tác giả để
khôi phục mạch trình bày, đồng thời chỉ giữ những kết luận phù hợp với cấu
trúc slide.

Bài toán gốc có ba học sinh \(A,B,C\). Mỗi người mang một số nguyên dương
trên trán, nhìn thấy hai số còn lại, và biết rằng trong ba số có một số bằng
tổng hai số kia. Giáo sư hỏi lần lượt từng người xem có đoán được số của
mình hay không. Sau một số câu trả lời ``không'', một người có thể dựa vào
chính các câu trả lời ấy để đoán ra.

\section{Vì sao cách nghĩ trực tiếp rườm rà}

Ví dụ kinh điển là ba số \(1,2,3\). Người \(A\) nhìn thấy \(2,3\), nên số của
mình có thể là \(1\) hoặc \(5\); \(A\) trả lời không. Người \(B\) nhìn thấy
\(1,3\), nên số của mình có thể là \(2\) hoặc \(4\); câu trả lời của \(A\)
chưa loại được trường hợp nào, nên \(B\) cũng trả lời không.

Đến \(C\), hai khả năng là \(1\) và \(3\). Nếu \(C\) giả sử số của mình là
\(1\), thì \(B\) đã nhìn thấy \(1,1\) và phải biết ngay số của mình là \(2\),
vì hiệu \(0\) không phải số nguyên dương. Nhưng thực tế \(B\) đã trả lời
không. Mâu thuẫn này loại bỏ khả năng \(1\), nên \(C\) đoán được số của mình
là \(3\).

Cách phân tích trên rất dễ hiểu, nhưng khi số lượt hỏi tăng, mỗi người phải
nghĩ xem người trước đang nghĩ gì về suy nghĩ của người trước nữa. Slide dùng
bài toán này để dẫn tới cách nhìn toàn cục: thay vì mô phỏng từng tầng suy
nghĩ, ta mô tả trạng thái và chuyển trạng thái.

\section{Trạng thái và ba loại tình huống}

Dùng bộ bốn
\[
  (a_1,a_2,a_3,k),\qquad k\in\{1,2,3\},
\]
để mô tả tình huống ba số là \(a_1,a_2,a_3\), quá trình hỏi bắt đầu từ học
sinh thứ nhất, và người đang được hỏi là học sinh thứ \(k\). Tập mọi bộ hợp
lệ được ký hiệu là \(M\).

Slide và bài viết phân biệt ba loại tình huống:
\begin{itemize}
  \item \term{tình huống kết thúc} \(W\): người đang được hỏi có thể trực tiếp
        xác định số của mình;
  \item \term{tình huống loại một} \(S_1\): số của người đang được hỏi bằng
        tổng hai số còn lại;
  \item \term{tình huống loại hai} \(S_2\): một trong hai người còn lại mang
        số bằng tổng.
\end{itemize}
Trong dạng gốc, nếu người \(k\) nhìn thấy hai số bằng nhau, khả năng hiệu
bằng \(0\) bị loại trực tiếp; đó là tình huống kết thúc.

Ký hiệu \(F(a_1,a_2,a_3,k)\) là số lượt hỏi sớm nhất để người đầu tiên đoán
ra trong tình huống tương ứng. Nếu chưa thể đoán, xem giá trị là vô hạn; trong
các tình huống của bài toán gốc, quá trình sẽ quy về kết thúc sau hữu hạn
bước.

\section{Truy hồi cho trường hợp ba người}

Giả sử đang ở tình huống loại một, tức là số của người \(k\) thật sự bằng
tổng hai số còn lại. Người đó muốn đoán đúng thì phải loại bỏ khả năng hiệu.
Việc loại bỏ khả năng hiệu tương đương với việc xét một tình huống mới, trong
đó số lớn nhất được thay bằng hiệu của hai số còn lại và lượt hỏi được chuyển
sang người đáng lẽ đã phát hiện mâu thuẫn.

Với \(k=1\), các công thức trong bài viết là:
\[
\begin{aligned}
F(a_1,a_2,a_3,1)&=1 &&(a_2=a_3),\\
F(a_1,a_2,a_3,1)&=2+F(a_2-a_3,a_2,a_3,2) &&(a_2>a_3),\\
F(a_1,a_2,a_3,1)&=1+F(a_3-a_2,a_2,a_3,3) &&(a_3>a_2).
\end{aligned}
\]
Hai người còn lại có công thức đối xứng:
\[
\begin{aligned}
F(a_1,a_2,a_3,2)&=2 &&(a_1=a_3),\\
F(a_1,a_2,a_3,2)&=1+F(a_1,a_1-a_3,a_3,1) &&(a_1>a_3),\\
F(a_1,a_2,a_3,2)&=2+F(a_1,a_3-a_1,a_3,3) &&(a_3>a_1),
\end{aligned}
\]
\[
\begin{aligned}
F(a_1,a_2,a_3,3)&=3 &&(a_1=a_2),\\
F(a_1,a_2,a_3,3)&=2+F(a_1,a_2,a_1-a_2,1) &&(a_1>a_2),\\
F(a_1,a_2,a_3,3)&=1+F(a_1,a_2,a_2-a_1,2) &&(a_2>a_1).
\end{aligned}
\]
Mỗi bước thay một số lớn bằng hiệu của hai số còn lại. Vì các số là nguyên
dương và giảm theo kiểu giống thuật toán Euclid, quá trình không thể kéo dài
vô hạn.

Kết luận trung tâm của slide là: trong dạng gốc, bất kể ba số cụ thể là gì,
người mang số lớn nhất sẽ là người đầu tiên đoán ra. Quan trọng hơn, ta đã
biến suy luận lồng nhau thành một truy hồi tuyến tính.

\section{Mở rộng nhiều người}

Slide tiếp tục mở rộng sang \(n\ge3\) học sinh. Khi một trong \(n\) số bằng
tổng \(n-1\) số còn lại, trạng thái có dạng
\[
  (a_1,a_2,\ldots,a_n,k).
\]
Với học sinh thứ \(k\), đặt
\[
  S_k=\sum_{i\ne k}a_i.
\]
Nếu \(v\ne k\) là chỉ số của số lớn nhất trong những số \(k\) nhìn thấy, hai
khả năng đáng xét cho số của \(k\) là
\[
  S_k \quad\text{và}\quad 2a_v-S_k.
\]
Nếu \(2a_v\le S_k\), khả năng thứ hai không dương và người đó rơi vào tình
huống kết thúc. Nếu \(2a_v>S_k\), ta chuyển sang trạng thái thay
\[
  a_k\leftarrow 2a_v-S_k
\]
và xét tiếp từ người \(v\). Độ lệch số lượt hỏi từ \(v\) tới \(k\) được tính
theo vòng tròn hỏi đáp. Truy hồi tổng quát vì vậy vẫn giữ tinh thần của
trường hợp ba người: loại bỏ khả năng sai bằng cách quy về một tình huống
nhỏ hơn.

\section{Mở rộng chia nhóm}

Mở rộng cuối trong bài viết, và cũng là phần mà slide có nhiều công thức
nhúng, là trường hợp giáo sư chia \(n\) học sinh thành hai nhóm có tổng bằng
nhau. Học sinh không biết cách chia nhóm, nên với người \(k\), mỗi cách chia
khả dĩ có thể sinh ra một giá trị khác nhau cho số trên trán mình.

Nếu nhóm đối diện \(k\) gồm các chỉ số \(u_1,\ldots,u_m\), một giá trị khả dĩ
có dạng
\[
  V_k=\sum_{r=1}^m a_{u_r}
      -\sum_{i\notin\{k,u_1,\ldots,u_m\}}a_i.
\]
Các giá trị không dương bị loại trực tiếp. Nếu còn đúng một giá trị dương,
đó là tình huống kết thúc; nếu còn nhiều giá trị dương, người \(k\) phải
chứng minh mọi giá trị sai đều dẫn tới mâu thuẫn với những câu trả lời trước
đó.

Khác với mở rộng ``một số bằng tổng các số còn lại'', ở đây các khả năng có
thể phân nhánh theo nhiều cách chia nhóm. Vì vậy quá trình suy luận không còn
là một đường duy nhất mà là một cây trạng thái. Tuy vậy, phương pháp vẫn
không đổi: mô tả tập khả năng, loại giá trị không dương, đệ quy trên từng
giá trị sai, và dùng ghi nhớ trạng thái khi lập trình.

\section{Thông điệp thuật toán}

Giá trị của bài trình bày không nằm ở một công thức riêng lẻ, mà ở cách bỏ
qua hình thức ``người này nghĩ người kia nghĩ gì'' để tìm cấu trúc bất biến.
Khi trạng thái được định nghĩa đúng, câu trả lời ``không'' của một người chỉ
là một phép loại trừ trong không gian trạng thái. Nhờ đó bài toán suy luận
tưởng như tăng theo cấp số mũ có thể được xử lý bằng truy hồi và quy hoạch
trạng thái.

\end{document}
